%% PArtie Analyse (Seif LAmouchi)

\chapter{Analyse et expression des besoins}

\section{Introduction générale}

L'évolution des réseaux informatiques et la généralisation des protocoles
de communication sécurisés ont considérablement modifié la nature du
trafic réseau. Des protocoles tels que HTTPS, TLS et QUIC assurent la
confidentialité des échanges, mais rendent également plus difficile
l'identification des applications à partir du contenu des paquets.

Dans ce contexte, la classification du trafic réseau constitue un enjeu
important pour les administrateurs réseau, notamment lorsqu'il s'agit
d'assurer une qualité de service (QoS) adaptée aux différentes
applications.

Ce projet porte sur la conception d'un agent hybride de classification
du trafic chiffré. Celui-ci combine des techniques d'apprentissage
automatique et des règles symboliques afin d'identifier les flux réseau,
de détecter les applications inconnues et de prendre des décisions
adaptées en matière de qualité de service.

Cette première partie présente le contexte du projet, la problématique,
les objectifs visés ainsi que les besoins fonctionnels et non
fonctionnels du système.

\section{Contexte du projet}

\subsection{Contexte général}

Les réseaux informatiques transportent simultanément plusieurs catégories
de trafic, correspondant à des applications ayant des exigences
différentes. On peut notamment distinguer :

\begin{itemize}
    \item La navigation Web, généralement tolérante à des variations
    modérées du délai.
    \item Le streaming vidéo, qui nécessite un débit suffisamment stable.
    \item La visioconférence et la voix sur IP (VoIP), sensibles à la
    latence et à la gigue.
    \item Les applications interactives, qui nécessitent une réponse
    rapide du réseau.
\end{itemize}

La coexistence de ces flux impose de disposer de mécanismes capables
d'identifier leur nature afin de leur appliquer des traitements réseau
appropriés.

Cependant, le chiffrement des communications limite l'accès aux
informations contenues dans les charges utiles des paquets. Les
mécanismes traditionnels de classification deviennent ainsi moins
efficaces lorsqu'ils reposent uniquement sur l'inspection du contenu.

\subsection{Problématique de la classification du trafic chiffré}

La classification traditionnelle peut s'appuyer sur les numéros de
ports ou sur l'inspection approfondie des paquets (Deep Packet
Inspection -- DPI).

Toutefois, ces approches présentent certaines limites :

\begin{itemize}
    \item Les numéros de ports ne permettent pas toujours d'identifier
    précisément une application.
    \item L'inspection approfondie est limitée lorsque les données
    transportées sont chiffrées.
    \item Les applications peuvent modifier leurs comportements ou
    utiliser des protocoles différents.
    \item Les méthodes classiques ne disposent pas nécessairement
    d'un mécanisme adapté à la reconnaissance des flux inconnus.
\end{itemize}

L'apprentissage automatique permet d'exploiter des caractéristiques
statistiques et temporelles des flux, sans nécessiter l'accès au contenu
déchiffré. Cependant, un modèle entraîné sur un ensemble de classes
connues peut attribuer à tort un flux inconnu à une classe existante \cite{drapergil2016,rezei2019}.

Il devient donc nécessaire de combiner plusieurs sources d'information
et de prévoir un mécanisme explicite de rejet des flux inconnus.

\subsection{Problématique du projet}

La problématique principale de ce projet peut être formulée comme suit :

\begin{quote}
Comment concevoir un agent hybride capable de classifier précocement
le trafic réseau chiffré, de reconnaître les flux inconnus et de
combiner apprentissage automatique et règles symboliques afin
d'améliorer la qualité de service des applications critiques ?
\end{quote}

Cette problématique soulève plusieurs défis :

\begin{itemize}
    \item \textbf{Classification précoce :} identifier un flux à partir
    des premiers paquets observés.
    \item \textbf{Reconnaissance en ensemble ouvert :} distinguer les
    applications connues des flux inconnus.
    \item \textbf{Fusion des connaissances :} combiner les prédictions
    du modèle appris avec les règles symboliques.
    \item \textbf{Robustesse :} évaluer le comportement du système face
    à la dérive des données et aux tentatives d'évasion.
    \item \textbf{Qualité de service :} appliquer une politique adaptée
    et mesurer son effet sur la latence et la gigue des flux critiques.
\end{itemize}

\section{Objectifs du projet}

L'objectif principal est de concevoir et d'évaluer un agent hybride
capable de classifier le trafic chiffré et d'exploiter les résultats
de cette classification pour piloter une politique de qualité de
service.

Les objectifs spécifiques sont les suivants :

\begin{enumerate}
    \item Extraire des caractéristiques statistiques et temporelles
    des flux réseau à partir de leurs premiers paquets.
    
    \item Développer et comparer un modèle de classification basé sur
    Random Forest et un modèle neuronal léger, tel qu'un CNN 1D.
    
    \item Mettre en place un mécanisme de rejet des flux inconnus,
    notamment à partir d'un seuil de confiance ou d'une détection
    hors distribution.
    
    \item Définir des règles symboliques indépendantes de
    l'apprentissage, exploitant notamment les ports, le SNI et les
    informations DNS lorsqu'elles sont disponibles.
    
    \item Concevoir un mécanisme d'arbitrage permettant de combiner
    les résultats du classifieur et ceux du moteur de règles.
    
    \item Appliquer une politique de qualité de service reposant
    notamment sur le marquage DSCP et la gestion des files HTB.
    
    \item Expliquer les décisions de classification et d'arbitrage
    à l'aide de SHAP.
    
    \item Évaluer les performances du système en mesurant notamment
    la précision de classification, le rejet des flux inconnus,
    le temps d'inférence, la latence et la gigue.
\end{enumerate}

L'évaluation reposera également sur la comparaison de trois
configurations : l'apprentissage seul, la méthode classique seule
et l'agent hybride.


\section{Étude de l'existant}

La classification du trafic réseau constitue une étape essentielle
dans la gestion et l'optimisation des réseaux informatiques. Plusieurs
approches ont été développées afin d'identifier les applications
générant les flux et de permettre leur traitement selon leurs besoins.

Les principales approches étudiées dans le cadre de ce projet sont
la classification basée sur les ports, l'inspection approfondie des
paquets, l'apprentissage automatique et les règles symboliques.

\subsection{Classification basée sur les numéros de ports}

La classification par ports est l'une des méthodes traditionnelles
d'identification du trafic réseau. Elle consiste à examiner les
numéros de ports source et destination associés aux protocoles de
transport, notamment TCP et UDP.

Certains services utilisent des ports conventionnels. Par exemple,
le port 80 est traditionnellement associé au protocole HTTP, tandis
que le port 443 est généralement utilisé par HTTPS.

Cette méthode présente l'avantage d'être simple à mettre en œuvre
et peu coûteuse en ressources. Cependant, elle devient moins fiable
lorsque plusieurs applications utilisent les mêmes ports ou lorsque
les applications emploient des ports dynamiques.

De plus, le numéro de port ne permet pas nécessairement d'identifier
précisément l'application à l'origine du trafic.

\textbf{Limite principale :} cette approche fournit un indice
d'identification, mais ne garantit pas une classification précise
des applications.

\subsection{Inspection approfondie des paquets (DPI)}

L'inspection approfondie des paquets, ou Deep Packet Inspection
(DPI), consiste à examiner différentes parties des paquets réseau,
notamment leurs en-têtes et leur contenu, afin d'identifier les
protocoles et les applications.

Contrairement à une classification reposant uniquement sur les
ports, le DPI peut exploiter des signatures et des motifs spécifiques
aux protocoles applicatifs.

Cette approche peut offrir une identification détaillée lorsque
les informations nécessaires sont accessibles. Toutefois, le
chiffrement des communications limite fortement l'inspection du
contenu applicatif.

Elle peut également nécessiter des ressources de calcul importantes,
en particulier lorsque le volume de trafic à analyser est élevé.

\textbf{Limite principale :} l'efficacité du DPI dépend de la
disponibilité des informations inspectables et peut être réduite
par le chiffrement des communications.

\subsection{Classification basée sur l'apprentissage automatique}

L'apprentissage automatique (Machine Learning) permet de classifier
le trafic réseau à partir de caractéristiques extraites des flux,
sans nécessiter systématiquement l'accès au contenu applicatif \cite{rezei2019}.

Ces caractéristiques peuvent notamment inclure :

\begin{itemize}
    \item La taille des paquets.
    \item La durée des flux.
    \item Les intervalles temporels entre les paquets.
    \item Le nombre de paquets échangés.
    \item La direction et la distribution des échanges.
\end{itemize}

Dans le cadre de ce projet, deux familles de modèles sont envisagées :

\begin{itemize}
    \item \textbf{Random Forest :} un modèle fondé sur un ensemble
    d'arbres de décision, adapté à la classification à partir de
    caractéristiques statistiques et temporelles \cite{breiman2001}.
    
    \item \textbf{CNN 1D :} un réseau neuronal convolutif capable
    d'exploiter des représentations séquentielles des caractéristiques
    des paquets.
\end{itemize}

L'apprentissage automatique permet ainsi d'identifier des comportements
de trafic plus complexes que ceux décrits par de simples règles
fixes.

Cependant, un modèle supervisé apprend à reconnaître les classes
présentes dans ses données d'entraînement. Face à une application
inconnue ou à une distribution de trafic différente, il peut
produire une prédiction incorrecte avec un niveau de confiance
potentiellement élevé \cite{scheirer2013,hendrycks2017}.

La classification précoce constitue également un défi : réduire
le nombre de paquets observés peut accélérer la décision, mais
risque de diminuer la précision.

\textbf{Limite principale :} un modèle appris seul ne garantit pas
la reconnaissance fiable des flux inconnus et peut être sensible
aux changements de distribution des données.

\subsection{Classification basée sur les règles symboliques}

Les règles symboliques reposent sur des conditions explicites
définies à partir de connaissances réseau. Elles permettent
d'associer certains indices observables à des catégories de trafic.

Dans le cadre de ce projet, les règles peuvent notamment exploiter :

\begin{itemize}
    \item Les numéros de ports source et destination.
    \item Les informations SNI, lorsqu'elles sont disponibles.
    \item Les informations DNS, lorsqu'elles sont accessibles.
    \item Certaines caractéristiques et conditions propres aux flux.
\end{itemize}

Contrairement à un modèle d'apprentissage automatique, les règles
symboliques sont définies explicitement et peuvent être examinées
directement par l'administrateur.

Elles présentent donc un avantage en matière d'interprétabilité
et permettent d'intégrer des connaissances réseau indépendantes
de l'entraînement d'un modèle.

Toutefois, leur efficacité dépend de la qualité et de la couverture
des règles définies. Une règle trop générale peut conduire à une
mauvaise classification, tandis qu'une base de règles limitée
peut ne pas couvrir les nouveaux comportements des applications.

\textbf{Limite principale :} les règles symboliques sont
interprétables, mais leur couverture dépend des connaissances
préalablement définies.

\subsection{Approche hybride : apprentissage automatique et règles}

L'approche hybride consiste à combiner les capacités de généralisation
d'un modèle appris avec les connaissances explicites représentées
par des règles symboliques \cite{garcez2023}. 

Dans le système étudié, les deux mécanismes fonctionnent comme
des sources de décision indépendantes :

\begin{itemize}
    \item Le classifieur fournit des probabilités ou des scores
    associés aux classes de trafic connues.
    
    \item Le moteur de règles analyse les indices réseau disponibles
    et produit des éléments de décision indépendants du modèle appris.
    
    \item Le module d'arbitrage combine ces informations afin de
    sélectionner une classe ou de rejeter explicitement le flux
    comme inconnu.
\end{itemize}

Cette architecture permet de ne pas dépendre exclusivement d'une
seule méthode de classification. Elle introduit également un
mécanisme de rejet afin d'éviter d'attribuer systématiquement
une classe connue à un flux qui ne peut pas être identifié de
manière suffisamment fiable.

Enfin, la classe retenue par l'arbitrage peut être utilisée par
le module d'action pour appliquer une politique de qualité de
service adaptée. Le module d'explication contribue à rendre
les décisions plus compréhensibles, notamment à l'aide de SHAP.

\subsection{Comparaison des approches}

Le tableau suivant synthétise les principales caractéristiques
des approches étudiées.

\begin{table}[H]
\centering
\caption{Comparaison des approches de classification du trafic}
\label{tab:comparaison-approches}

\small
\begin{tabularx}{\textwidth}{|p{2.5cm}|Y|Y|Y|}
\hline
\textbf{Approche} &
\textbf{Principe} &
\textbf{Avantages} &
\textbf{Limites} \\
\hline

Classification par ports &
Analyse des ports source et destination. &
Simple, rapide et peu coûteuse. &
Identification imprécise lorsque plusieurs applications partagent
les mêmes ports. \\
\hline

DPI &
Inspection des en-têtes et du contenu des paquets. &
Identification détaillée lorsque le contenu est accessible. &
Limitée par le chiffrement et potentiellement coûteuse en ressources. \\
\hline

Machine Learning &
Classification à partir de caractéristiques statistiques et temporelles. &
Capacité à apprendre des comportements complexes. &
Sensibilité aux flux inconnus et aux changements de distribution. \\
\hline

Règles symboliques &
Application de conditions définies à partir de connaissances réseau. &
Interprétabilité et indépendance vis-à-vis de l'entraînement. &
Couverture limitée par les règles disponibles. \\
\hline

Approche hybride &
Combinaison du modèle appris, des règles et d'un mécanisme d'arbitrage. &
Diversification des sources de décision et possibilité de rejet des
flux inconnus. &
Complexité supplémentaire liée à l'arbitrage et à la coordination
des modules. \\
\hline

\end{tabularx}
\end{table}

\subsection{Justification du choix de l'approche hybride}

L'étude des différentes approches montre qu'aucune méthode ne répond
à elle seule à l'ensemble des contraintes du projet.

La classification par ports et le DPI présentent des limites
face à l'évolution des protocoles et au chiffrement. L'apprentissage
automatique permet d'exploiter des caractéristiques comportementales,
mais nécessite un mécanisme adapté à la détection des flux inconnus.
Les règles symboliques apportent des connaissances explicites,
mais ne peuvent couvrir à elles seules tous les comportements
possibles.

Le choix d'une architecture hybride permet donc de combiner ces
différentes sources d'information au sein d'un même système.

Cette approche constitue la base de conception retenue pour le
projet. Son intérêt devra être vérifié expérimentalement en
comparant trois configurations : apprentissage automatique seul,
méthode classique seule et agent hybride.

